\section{R3: Evaluation of programs with holes}\label{sec:r3}

\question{R3.Q1}{What is already available for higher-order and dependent examples?}
\status{Higher-order examples have a concrete precedent in Smyth; dependent closures need additional typed-substitution machinery.}

Smyth represents incomplete computations by hole closures and propagates
required results backward into constraints on holes. Its implementation
supports higher-order function examples even though its core exposition is
simplified to first-order examples \cite{smyth}. Leant2 should begin by
adapting that distinction: higher-order behavior is observed through a finite
collection of applications, not by demanding a complete mathematical function
as the output of unevaluation.

For a hole $h$ declared under telescope $\Delta$, represent an observation as
\[
 (h,\rho,\mathcal O),\qquad
 \rho:\Gamma\longrightarrow\Delta,
 \tag{6}\label{eq:closure}
\]
where $\rho$ is a well-typed environment substitution and $\mathcal O$ is a
typed elimination or result constraint. Two occurrences of $h$ can have
substitutions that instantiate indices differently. A hole producing
$\mathsf{Vec}\ B\ n$ may be observed at length zero in one closure and length
three in another. These are constraints on one contextual term, not two
untyped examples of a common ground output type.

The first dependent difficulty is therefore not merely storing proofs in an
environment. It is transporting the expected observation type along $\rho$
and composing substitutions correctly under binders. The next difficulty is
backward propagation through an eliminator whose motive depends on the value
being inspected. Begin with a typed first-order data fragment, then dependent
constructors with known indices, then application observations of function
holes. Unsupported inversions should produce residual constraints rather than
inventing an untyped example.

A contract is not necessarily a list of equations. Normalize its Boolean and
propositional structure using checked equivalences, retaining provenance to
the original proposition. Splitting \code{List.all} over a known finite list
or a conjunction of Booleans is useful. Replacing \code{a == b = true} by
$a=b$ needs a suitable lawfulness theorem; \code{BEq} alone does not supply
one. Multiple calls to the synthesized function in one observation must remain
related rather than being treated as independent input-output examples.

No source examined here supplies a ready-made, conservative evaluator for the
whole Lean language with this dependent closure interface. That identifies
an implementation and proof task, not evidence that every component of it is
unexplored.

\question{R3.Q2}{How should normalization caches survive assignment and rollback?}
\status{Use dependency-tracked persistent memoization; do not key validity only by hole identifiers or a global update counter.}

Self-adjusting computation provides a general precedent for memoization with
change propagation and a consistency theorem \cite{selfadjust}. It is not a
drop-in normalizer for Lean. The adaptation is to record what an evaluation
\emph{read}, not only what variable it was visibly stuck on.

A memo entry should contain the expression or closure identity, local-context
fingerprint, relevant environment and transparency data, a result, and a read
set. Each read-set item records an assignment's identity or immutable version.
Include universe assignments, delayed-assignment links, instance dictionaries,
and type dependencies whenever the evaluator inspected them. Reuse is valid
only when all relevant reads still denote the same values.

\begin{proposition}[Read-set cache reuse]\label{prop:cache}
Assume a deterministic evaluator's result depends only on an immutable
expression, fixed semantic parameters, and the values in its recorded read
set. If another state agrees on all those values, the memoized result is valid
there.
\end{proposition}
\begin{proof}
Replay the evaluation. Every read yields the recorded value, so every branch
and reduction rule is the same. Induction over the finite evaluation trace
gives the same result. The completeness of the read set is an explicit
hypothesis, not a property implied by recording just blocked holes.
\end{proof}

A rollback can restore an old assignment and later take a different branch
using the same metavariable identifier. This is the familiar identity-reuse
problem: the name did not change, but its meaning did. Immutable state tokens,
content identities, or validated assignment fingerprints avoid it. A single
global timestamp is safe if it invalidates everything on every relevant change,
but loses most reuse; a reused timestamp is unsafe.

Use per-observation caches first. The current code already decomposes
conjunctions and remembers blockers, so this is an incremental improvement.
However, \code{residualBlockers} is outside the saved meta-state. A future
frontier must attach its memo identity to the node, or validate the entire
read set before reuse \cite{searchcode,transactioncode}. False results can be
cached globally only with a proof or a key strong enough to establish that
the exact same proposition and program are being refuted.

\question{R3.Q3}{What is the cheapest route to a provably conservative fast evaluator?}
\status{A small proof-producing or reflected fragment is preferable to a second unverified implementation of the entire kernel.}

The evaluator need not decide every Lean term. It needs to prove the decisions
it does make. Define a typed observation language containing the constructors,
projections, primitive operations, and recursors the synthesizer actually
uses. Give it a denotation into Lean terms and prove a fuelled evaluation
correctness theorem. A successful evaluation can then provide a theorem about
the original observation; unsupported syntax or fuel exhaustion gives
\emph{unknown}.

There are two practical implementation paths. A proof-producing reducer emits
an equality certificate for each step, composing certificates only for the
parts inspected. A reflected evaluator runs over a reified typed syntax and
uses one proved soundness theorem. The former has simpler incremental
integration; the latter may give smaller checking overhead after a larger
initial proof investment. Both permit a fast untrusted proposal phase followed
by validation of the claimed reduction.

\begin{proposition}[Residual refutation]\label{prop:eval}
Suppose a partial program $p$ has a well-typed completion telescope $\Delta$.
If a checked certificate proves
\[
 \forall\theta:\Delta,\quad
 \mathsf{decide}(C(p[\theta]))=\mathsf{false},
\]
using an appropriate decision procedure for each instance, then rejecting $p$
is conservative. A certificate under additional assumptions is sufficient only
when those assumptions follow from the requested contract for every accepted
completion.
\end{proposition}
\begin{proof}
For any completion, instantiate the certificate. The correctness of
\code{decide} yields $\neg C(p[\theta])$. Under additional assumptions, a
hypothetical accepted completion supplies those assumptions and produces the
same contradiction.
\end{proof}

The current kernel-reduction path has a strong advantage: genuine reduction
of a well-typed expression commutes with well-typed substitution. But its
\code{instantiatePartial} helper explicitly constructs reduction-only
expressions whose metavariable contexts may no longer match. A completeness
argument still needs a well-scoped closure/substitution account for that
transformation. Calling the kernel's reducer does not by itself prove a
separate preprocessing transformation correct \cite{searchcode}.

Lean4Lean is a useful source of formalized kernel structure and verification
work, not a button that extracts an incremental closure evaluator
\cite{lean4lean}. A closure machine has its own environment relation and
substitution lemmas. Reuse relevant definitions and proved lemmas where they
match, but prove the fragment's simulation theorem explicitly. Differential
testing remains valuable for finding mistakes; even a hundred thousand passing
closed tests is not a proof about every open completion.

\question{R3.Q4}{Can top-down angelic recursion have termination and completeness guarantees?}
\status{Yes for carefully bounded fragments with fair exploration and relational oracle constraints; not by simply guessing one helpful return value per call.}

An angelic recursive call denotes an existentially chosen result subject to
constraints. Equal arguments to the same pure function must share a result.
For example, independently choosing two results could satisfy
$f(0)\ne f(0)$, although no function can. Maintain a map of call observations
with equality and relational constraints, rather than unrelated successful
oracle guesses.

A supplied equation $f(v)=w$ can constrain the oracle at $v$ because every
accepted function must respect it. An unobserved recursive result is different:
a convenient value may guide a branch, but failure of that value does not
refute all possible results. To prune using angelic execution, compute an
over-approximation of possible results and prove it cannot satisfy the
observation. Otherwise use angelic information only for ordering.

Burst gives a concrete refinement-based synthesis design, but its arguments
cannot be transferred merely by changing bottom-up enumeration to top-down
search \cite{burst}. A top-down bounded theorem needs finite skeletons, finite
oracle domains or a terminating complete constraint solver, fair exploration
of refinements, and a decreasing or bounded measure preventing repeated
unchanged refinements. Under those conditions the search terminates; if its
rules cover every bounded candidate and all pruning is conservative, it is
complete for that fragment. For unbounded syntax, dovetailing yields at most
a relative semidecision guarantee. No general termination claim follows from
the existence of a final kernel gate.
